\documentclass[10pt,a4paper]{amsart}
\usepackage[margin=19mm]{geometry}
\usepackage{amsmath,amssymb,mathtools}
\usepackage[scr=rsfs,cal=euler]{mathalfa}
\usepackage{xfrac}
\usepackage[unicode,hidelinks,bookmarks=false]{hyperref}
\newcommand{\ie}{i.e.\ }
\newcommand{\cf}{cf.\ }
\newcommand{\resp}{respectively\ }
\newcommand{\Subsection}{\subsection}
\providecommand{\nobreakdash}{\nobreak\hbox{-}}
\setlength{\emergencystretch}{2em}
\multlinegap0pt
\usepackage{bm}
\usepackage{bbm}
\usepackage{dsfont}
\usepackage{xspace}
\usepackage{cancel}
\usepackage{mathtools}
\usepackage{amsthm}
\makeatletter
\@ifundefined{theorem}{\newtheorem{theorem}{Theorem}}{}
\@ifundefined{lemma}{\newtheorem{lemma}[theorem]{Lemma}}{}
\makeatother
\title[Characterizations of bipartite and Eulerian partial duals of]{Characterizations of bipartite and Eulerian partial duals of orientable hypermaps}
\author{Yufan Han, Metrose Metsidik}
\date{Source: arXiv:2606.30071v1 (CC BY 4.0)}
\begin{document}
\maketitle

\noindent\emph{Source and licence.} Characterizations of bipartite and Eulerian partial duals of orientable hypermaps. Version arXiv:2606.30071v1 (CC BY 4.0), distributed under the Creative Commons Attribution 4.0 International licence, as recorded in the arXiv licence metadata for this exact version and shown on the abstract page https://arxiv.org/abs/2606.30071v1. Full rights notice: licenses/arxiv-2606.30071-CC-BY-4.0.txt.\par\smallskip

\noindent\textit{Editorial scope.} Counted unit: the Lemma whose source label is \texttt{lem:checkerboard} in the article (source0.tex lines 329--338, source label \texttt{lem:checkerboard}), delivered together with the article's own complete original proof of it (source0.tex lines 340--446, one \texttt{proof} environment of 107 lines). No mathematics is written, repaired or re-derived: statement and proof are the article's own text line for line. Required context retained uncounted: none. Every definition, display and label of those lines is retained unchanged. Omitted from this package and not redistributed: 1--328, 339--339, 447--1686. Nothing inside a retained excerpt is omitted or altered: every retained body is delivered exactly as the article's lines are written, so no omission receipt of any kind accompanies this package. Citations: the retained excerpts carry no citation command at all, so no bibliography entry of the article is transcribed and no reference is redistributed. Reference adaptation: none was needed and none was made; every reference inside the delivered unit resolves inside it, so no unresolved reference or citation is produced. Printed slips kept literal: no printed word, symbol, hypothesis or number of the counted statement or of its proof was altered. Layout and class: the article's own class is \texttt{article} with packages including arxiv, inputenc, fontenc, hyperref, url, booktabs, amsfonts, nicefrac, microtype, lipsum, graphicx, amsmath, amssymb, amsthm; the wrapper is the lane's pinned \texttt{amsart} editorial wrapper at 10pt on A4, which restates the theorem-like environments the retained text needs and adds the article's own macro definitions that the retained text needs (none), with extra wrapper packages (bm, bbm, dsfont, xspace, cancel, mathtools). The delivered numbering is the unit's own. Assets: no retained span contains a figure environment, a graphics-inclusion command, a PDF-page inclusion, a table or a TikZ picture; no image file is copied into this package, the package declares no asset dependency and no image file is redistributed with it. Licence: version arXiv:2606.30071v1 of this article is distributed under the Creative Commons Attribution 4.0 International licence, as recorded in the arXiv licence metadata for this exact version and shown on the abstract page https://arxiv.org/abs/2606.30071v1. A full rights notice is delivered at licenses/arxiv-2606.30071-CC-BY-4.0.txt. Characters: every retained excerpt is pure ASCII, so the delivered engine, XeLaTeX with its default Latin Modern text font, reports no missing character.
\begin{lemma}\label{lem:checkerboard}
Let $\phi_M:=\alpha_M^{-1}\sigma_M=\alpha_M\sigma_M$ denote the face permutation of the medial map $M(H)$ associated with hypermap $H=(\sigma,\alpha)$. Then:
\begin{enumerate}
  \item the restriction of $\phi_M$ to the negative points $\{i^-:i\in D\}$ is isomorphic to $\sigma$;
  \item the restriction of $\phi_M$ to the positive points $\{i^+:i\in D\}$ is isomorphic to $\sigma^{-1}\alpha$, and hence has the same cycle decomposition as the original face permutation $\phi=\alpha^{-1}\sigma$;
  \item the medial map \(M(H)\) admits a canonical checkerboard coloring: the
\(\phi_M\)-cycles contained in \(D^-\) are colored black, and the
\(\phi_M\)-cycles contained in \(D^+\) are colored white. 
\end{enumerate}
\end{lemma}

\begin{proof}
First take a negative point \(i^-\). By the definition of \(\sigma_M\), we have
\[
    \sigma_M(i^-)=i^+.
\]
Since \(\alpha_M\) contains the transposition \((i^+\,\sigma(i)^-)\), it follows that
\[
    \phi_M(i^-)
    =\alpha_M\sigma_M(i^-)
    =\alpha_M(i^+)
    =\sigma(i)^-.
\]
Thus the restriction of \(\phi_M\) to \(D^-\) is identified with \(\sigma\).

Now take a positive point \(i^+\). Write \(\alpha(i)=j\). By the construction of
\(\sigma_M\), the point following \(i^+\) around the corresponding medial vertex is
\(j^-\). Hence
\[
    \sigma_M(i^+)=j^-=\alpha(i)^-.
\]
Moreover, since \(\alpha_M\) contains the transposition
\[
    \bigl(\sigma^{-1}(j)^+\,j^-\bigr),
\]
we obtain
\[
    \phi_M(i^+)
    =\alpha_M\sigma_M(i^+)
    =\alpha_M(j^-)
    =\sigma^{-1}(j)^+
    =\sigma^{-1}\alpha(i)^+.
\]
Therefore the restriction of \(\phi_M\) to \(D^+\) is identified with
\(\sigma^{-1}\alpha\). Since
\[
    \sigma^{-1}\alpha=(\alpha^{-1}\sigma)^{-1}=\varphi^{-1},
\]
this restriction has the same cycle decomposition as the original face permutation
\(\phi=\alpha^{-1}\sigma\), with the cyclic order of each cycle reversed.

By parts (1) and (2), the face permutation
\[
        \phi_M=\alpha_M^{-1}\sigma_M=\alpha_M\sigma_M
\]
satisfies
\[
        \phi_M(D^-)\subseteq D^-,
        \qquad
        \phi_M(D^+)\subseteq D^+ .
\]
Hence every face-cycle of \(M(H)\) is contained entirely in either
\(D^-\) or \(D^+\). We color the face-cycles contained in \(D^-\) black
and the face-cycles contained in \(D^+\) white.

We now describe this coloring locally at a vertex of the medial map. Let
\[
        e=(i_1\,i_2\,\cdots\,i_k)
\]
be a hyperedge of \(H\). By the definition of the medial map, the
corresponding vertex \(v_e\) of \(M(H)\) has cyclic order
\[
        i_1^-\, i_1^+\, i_2^-\, i_2^+\,\cdots\, i_k^-\, i_k^+ .
\]
Thus the corners at \(v_e\) are of two alternating forms:
\[
        (i_j^-,i_j^+)
        \qquad\text{and}\qquad
        (i_j^+,i_{j+1}^-),
        \quad i_{k+1}:=i_1 .
\]

We claim that the corners of the first form are black corners and the
corners of the second form are white corners. Indeed, for a negative
point \(i_j^-\), the face traversal gives
\[
        i_j^-
        \xrightarrow{\ \sigma_M\ }
        i_j^+
        \xrightarrow{\ \alpha_M\ }
        \sigma(i_j)^- .
\]
Thus, when the corresponding face passes through \(v_e\), it passes
through the local corner \((i_j^-,i_j^+)\). Since the corresponding
\(\phi_M\)-cycle is contained in \(D^-\), this is a black face-corner.

On the other hand, for a positive point \(i_j^+\), we have
\[
        i_j^+
        \xrightarrow{\ \sigma_M\ }
        i_{j+1}^-
        \xrightarrow{\ \alpha_M\ }
        \sigma^{-1}(i_{j+1})^+
        =
        \sigma^{-1}\alpha(i_j)^+ .
\]
Thus, when the corresponding face passes through \(v_e\), it passes
through the local corner
\[
        (i_j^+,i_{j+1}^-)
        =
        (i_j^+,\alpha(i_j)^-).
\]
Since the corresponding \(\phi_M\)-cycle is contained in \(D^+\), this is
a white face-corner. 

Thus black and white face corners alternate at every vertex of \(M(H)\). Accordingly, such black-and-white coloring yields a checkerboard coloring of \(M(H)\).
\end{proof}

\end{document}
